% !TeX root = ../../main.tex
% ============================================================================
% MODULE 8: Analysis and Verification of Embedded Systems
% EE322 — Embedded Systems Design
% ============================================================================

\chapter{Analysis and Verification of Embedded Systems}
\label{ch:analysis}
\chaptermark{Module 8}

\begin{moduleinfo}
  \textbf{Module:} 8 \quad
  \textbf{Lecture Hours:} 4 \quad
  \textbf{ILOs Addressed:} ILO-6, ILO-7 \\[4pt]
  \textbf{Learning Outcomes:} By the end of this module, students will be able to:
  \begin{enumerate}[nosep]
    \item Profile embedded software for execution time, memory usage, and power consumption.
    \item Perform Failure Mode and Effects Analysis (FMEA) on embedded systems.
    \item Calculate reliability metrics (MTTF, MTBF, availability, FIT rate).
    \item Apply formal verification concepts (model checking with SPIN, temporal logic).
    \item Explain MISRA-C guidelines and their role in safety-critical software.
  \end{enumerate}
\end{moduleinfo}

% ----------------------------------------------------------------------------
\section{Performance Analysis and Optimisation}
\label{sec:ana:perf}

\subsection{Profiling Metrics}

Key performance indicators for embedded systems:
\begin{itemize}[nosep]
  \item \textbf{Execution time:} CPU cycles or wall-clock time consumed by a function or task.
  \item \textbf{Memory footprint:} Flash (code + constants) and RAM (stack + heap + global data).
  \item \textbf{Power consumption:} Average and peak current draw; energy per operation.
  \item \textbf{Throughput:} Data processed per unit time (bytes/s, samples/s).
  \item \textbf{Latency:} Time from stimulus to response (interrupt latency, task response time).
\end{itemize}

\subsection{ARM Cortex-M DWT Cycle Counter}

The Data Watchpoint and Trace (DWT) unit on Cortex-M3/M4/M7 provides a free-running cycle counter.

\begin{lstlisting}[style=C, caption={Using the DWT cycle counter for execution time measurement.}]
#include <stdint.h>

// DWT register addresses
#define DWT_CYCCNT   (*(volatile uint32_t *)0xE0001004)
#define DWT_CONTROL  (*(volatile uint32_t *)0xE0001000)
#define SCB_DEMCR    (*(volatile uint32_t *)0xE000EDFC)

void dwt_init(void) {
    SCB_DEMCR |= (1 << 24);    // Enable DWT
    DWT_CYCCNT = 0;             // Reset counter
    DWT_CONTROL |= 1;           // Enable CYCCNT
}

static inline uint32_t dwt_get_cycles(void) {
    return DWT_CYCCNT;
}

// Usage example
void measure_function(void) {
    uint32_t start = dwt_get_cycles();
    
    // --- Code under measurement ---
    my_dsp_function();
    // ---
    
    uint32_t end = dwt_get_cycles();
    uint32_t elapsed = end - start;  // handles wrap-around
    
    // At 168 MHz: time_us = elapsed / 168
    float time_us = (float)elapsed / 168.0f;
}
\end{lstlisting}

\subsection{Memory Analysis}

\begin{lstlisting}[style=C, caption={Checking flash and RAM usage from the linker map file.}]
/* After building, use arm-none-eabi-size: */
//    text    data     bss     dec     hex  filename
//   24576    1024    8192   33792    8400  firmware.elf
//
// Flash usage: text + data = 24576 + 1024 = 25600 bytes
// RAM usage:   data + bss  = 1024 + 8192  = 9216 bytes
\end{lstlisting}

\subsection{Optimisation Techniques}

\begin{table}[H]
\centering
\caption{Common embedded optimisation strategies.}
\label{tab:ana:optim}
\begin{tabularx}{\linewidth}{@{}lX@{}}
\toprule
\textbf{Technique} & \textbf{Description} \\
\midrule
Compiler optimisation flags & \texttt{-O2} (speed), \texttt{-Os} (size), \texttt{-Og} (debug-friendly) \\
Loop unrolling & Reduce loop overhead at the cost of code size \\
Lookup tables & Replace computation with pre-computed table (e.g., sine LUT) \\
Fixed-point arithmetic & Avoid floating-point on cores without FPU \\
DMA transfers & Offload data movement from CPU \\
Sleep modes & Reduce power by entering low-power states when idle \\
Dead code elimination & Remove unused functions (\texttt{-ffunction-sections -Wl,--gc-sections}) \\
\bottomrule
\end{tabularx}
\end{table}

\begin{practicalNote}
Always profile before optimising. Premature optimisation is the root of all evil (Knuth). Use the DWT counter or an oscilloscope to identify hotspots.
\end{practicalNote}

% ----------------------------------------------------------------------------
\section{Failure Mode and Effects Analysis (FMEA)}
\label{sec:ana:fmea}

\begin{definition}[FMEA]
Failure Mode and Effects Analysis is a systematic, proactive method for analysing potential failure modes of a system, their causes, effects, and criticality. The output is a prioritised list of risks, each assigned a Risk Priority Number (RPN).
\end{definition}

\[
\text{RPN} = \text{Severity} \times \text{Occurrence} \times \text{Detection}
\]

where each factor is rated 1--10 (1 = best, 10 = worst).

\begin{table}[H]
\centering
\caption{Example FMEA table for an embedded temperature controller.}
\label{tab:ana:fmea}
\footnotesize
\begin{tabularx}{\linewidth}{@{}p{1.8cm}p{2.2cm}Xp{0.6cm}p{0.6cm}p{0.6cm}p{0.8cm}X@{}}
\toprule
\textbf{Component} & \textbf{Failure Mode} & \textbf{Effect} & \textbf{S} & \textbf{O} & \textbf{D} & \textbf{RPN} & \textbf{Mitigation} \\
\midrule
ADC & Stuck at max & Controller thinks temp is max, turns heater off permanently & 7 & 3 & 5 & 105 & Range check on ADC value \\
MCU & Watchdog timeout & System resets, brief control loss & 5 & 2 & 2 & 20 & Safe state on reset \\
Heater relay & Stuck closed & Uncontrolled heating, fire risk & 10 & 2 & 6 & 120 & Thermal fuse, redundant temp sensor \\
Firmware & Stack overflow & Undefined behaviour, crash & 8 & 4 & 7 & 224 & Stack canary, FreeRTOS stack check \\
Power supply & Brownout & MCU resets or corrupts data & 6 & 3 & 4 & 72 & BOR circuit, EEPROM write verify \\
\bottomrule
\end{tabularx}
\end{table}

% ----------------------------------------------------------------------------
\section{Reliability and Availability}
\label{sec:ana:reliability}

\begin{definition}[Reliability]
Reliability $R(t)$ is the probability that a system performs its intended function without failure for a specified time $t$ under stated conditions. For constant failure rate $\lambda$:
\[
R(t) = e^{-\lambda t}
\]
\end{definition}

Key metrics:
\begin{align}
\text{MTTF} &= \frac{1}{\lambda} & &\text{(Mean Time To Failure, non-repairable)} \\
\text{MTBF} &= \text{MTTF} + \text{MTTR} & &\text{(Mean Time Between Failures)} \\
\text{Availability} &= \frac{\text{MTTF}}{\text{MTTF} + \text{MTTR}} & &\text{(Steady-state)} \\
\text{FIT rate} &= \frac{10^9}{\text{MTTF (hours)}} & &\text{(Failures In Time per $10^9$ hours)}
\end{align}

\subsection{Bathtub Curve}

\begin{figure}[H]
\centering
\begin{tikzpicture}
  \begin{axis}[
    width=12cm, height=6cm,
    xlabel={Time},
    ylabel={Failure Rate $\lambda(t)$},
    xmin=0, xmax=10, ymin=0, ymax=5,
    xtick=\empty, ytick=\empty,
    axis lines=left,
    every axis x label/.style={at={(axis description cs:1,0)}, anchor=west},
    every axis y label/.style={at={(axis description cs:0,1)}, anchor=south},
    clip=false
  ]
    % Infant mortality (decreasing)
    \addplot[domain=0:2.5, thick, ee-blue, samples=50] {3*exp(-1.5*x) + 0.5};
    
    % Useful life (constant)
    \addplot[domain=2.5:7, thick, ee-green!70!black, samples=50] {0.5};
    
    % Wear-out (increasing)
    \addplot[domain=7:10, thick, ee-red, samples=50] {0.5 + 0.15*(x-7)^2};
    
    % Annotations
    \node[font=\footnotesize, text=ee-blue] at (axis cs:1.2,3.5) {Infant Mortality};
    \node[font=\footnotesize, text=ee-green!70!black] at (axis cs:4.75,1.2) {Useful Life};
    \node[font=\footnotesize, text=ee-green!70!black] at (axis cs:4.75,0.9) {(constant $\lambda$)};
    \node[font=\footnotesize, text=ee-red] at (axis cs:8.8,2.5) {Wear-Out};
    
    % Dashed dividers
    \draw[dashed, thin, gray] (axis cs:2.5, 0) -- (axis cs:2.5, 4.5);
    \draw[dashed, thin, gray] (axis cs:7, 0) -- (axis cs:7, 4.5);
    
  \end{axis}
\end{tikzpicture}
\caption{Bathtub curve: failure rate over the lifetime of electronic systems.}
\label{fig:ana:bathtub}
\end{figure}

\subsection{Redundancy: Triple Modular Redundancy (TMR)}

\begin{figure}[H]
\centering
\begin{tikzpicture}[
    block/.style={draw, thick, minimum width=2cm, minimum height=1cm, rounded corners=2pt, font=\small\bfseries},
    arrow/.style={-{Stealth[length=3mm]}, thick}
  ]
  
  \node[block, fill=ee-blue!15] (m1) at (0, 1.5) {Module 1};
  \node[block, fill=ee-blue!15] (m2) at (0, 0)   {Module 2};
  \node[block, fill=ee-blue!15] (m3) at (0,-1.5)  {Module 3};
  
  \node[block, fill=ee-amber!15, minimum width=2.5cm] (voter) at (5, 0) {Majority Voter};
  
  \node[font=\small] (input) at (-3.5, 0) {Input};
  \node[font=\small] (output) at (8.5, 0) {Output};
  
  \draw[arrow] (input) -- (-1.5,0);
  \draw[arrow] (-1.5,0) |- (m1);
  \draw[arrow] (-1.5,0) -- (m2);
  \draw[arrow] (-1.5,0) |- (m3);
  
  \draw[arrow] (m1) -| (voter);
  \draw[arrow] (m2) -- (voter);
  \draw[arrow] (m3) -| (voter);
  
  \draw[arrow] (voter) -- (output);
  
\end{tikzpicture}
\caption{Triple Modular Redundancy: three identical modules feed a majority voter.}
\label{fig:ana:tmr}
\end{figure}

TMR reliability (assuming independent failures, perfect voter):
\[
R_{\text{TMR}}(t) = 3R(t)^2 - 2R(t)^3 = 3e^{-2\lambda t} - 2e^{-3\lambda t}
\]

TMR is more reliable than a single module for $R(t) > 0.5$ (i.e., early in the system's life).

% ----------------------------------------------------------------------------
\section{Formal Verification}
\label{sec:ana:formal}

\begin{definition}[Formal Verification]
Formal verification uses mathematical techniques to prove or disprove that a system design satisfies a formal specification. Unlike testing, which can only show the presence of bugs, formal verification can prove their \emph{absence} (within the scope of the model).
\end{definition}

\subsection{Model Checking with SPIN and Promela}

SPIN (Simple Promela Interpreter) is a model checker for verifying properties of concurrent systems.

\begin{lstlisting}[language=C, style=C, caption={Promela model of a producer-consumer system with bounded buffer.}, label={lst:ana:promela}]
/* Promela: producer-consumer with semaphores */
#define N 4  /* buffer size */

byte buffer[N];
byte in = 0, out = 0, count = 0;

/* Semaphores modeled as byte variables */
byte mutex = 1;
byte empty_slots = N;
byte full_slots = 0;

inline P(sem) { atomic { sem > 0 -> sem-- } }
inline V(sem) { atomic { sem++ } }

active proctype Producer() {
    byte item = 0;
    do
    :: P(empty_slots);
       P(mutex);
       buffer[in] = item;
       in = (in + 1) % N;
       count++;
       V(mutex);
       V(full_slots);
       item++;
    od
}

active proctype Consumer() {
    byte data;
    do
    :: P(full_slots);
       P(mutex);
       data = buffer[out];
       out = (out + 1) % N;
       count--;
       V(mutex);
       V(empty_slots);
    od
}

/* Safety property: buffer never overflows or underflows */
ltl safe { [](count >= 0 && count <= N) }
\end{lstlisting}

\subsection{Temporal Logic}

Temporal logic expresses properties about system behaviour over time:

\textbf{Linear Temporal Logic (LTL):}
\begin{itemize}[nosep]
  \item $\square p$ (``always $p$''): $p$ holds in every state.
  \item $\lozenge p$ (``eventually $p$''): $p$ holds in some future state.
  \item $p \mathbin{\mathcal{U}} q$ (``$p$ until $q$''): $p$ holds until $q$ becomes true.
  \item $\bigcirc p$ (``next $p$''): $p$ holds in the next state.
\end{itemize}

Common embedded system properties expressed in LTL:
\begin{itemize}[nosep]
  \item \textbf{Safety} (bad things never happen): $\square \neg \text{overflow}$
  \item \textbf{Liveness} (good things eventually happen): $\square(\text{request} \rightarrow \lozenge \text{grant})$
  \item \textbf{Mutual exclusion:} $\square \neg (\text{cs}_1 \land \text{cs}_2)$
  \item \textbf{Deadlock freedom:} $\square \lozenge \text{progress}$
\end{itemize}

\subsection{Timed Automata and UPPAAL}

UPPAAL models real-time systems as networks of timed automata with clock constraints. Useful for verifying timing properties like ``the airbag deploys within 10\,ms of impact detection.''

% ----------------------------------------------------------------------------
\section{MISRA-C Guidelines}
\label{sec:ana:misra}

\begin{definition}[MISRA-C]
MISRA-C\cite{misrac2012} is a set of software development guidelines for the C language in safety-critical and security-critical embedded systems. Originally developed for the automotive industry, it is now used across aerospace, medical, rail, and defence.
\end{definition}

\textbf{MISRA-C 2012 rule categories:}
\begin{itemize}[nosep]
  \item \textbf{Mandatory:} Must be followed (e.g., no undefined behaviour).
  \item \textbf{Required:} Must be followed unless a formal deviation is documented.
  \item \textbf{Advisory:} Recommended best practice.
\end{itemize}

\textbf{Selected important rules:}

\begin{table}[H]
\centering
\caption{Selected MISRA-C 2012 rules relevant to embedded programming.}
\label{tab:ana:misra}
\begin{tabularx}{\linewidth}{@{}lp{2.2cm}X@{}}
\toprule
\textbf{Rule} & \textbf{Category} & \textbf{Description} \\
\midrule
Rule 1.3 & Required & There shall be no occurrence of undefined or critical unspecified behaviour. \\
Rule 11.3 & Required & A cast shall not be performed between a pointer to object type and a pointer to a different object type. \\
Rule 14.3 & Required & Controlling expressions shall not be invariant. \\
Rule 17.7 & Required & The value returned by a function having non-void return type shall be used. \\
Rule 21.3 & Required & The memory allocation and deallocation functions of \texttt{<stdlib.h>} shall not be used. \\
Dir 4.1 & Required & Run-time failures shall be minimised. \\
Dir 4.6 & Advisory & \texttt{typedefs} that indicate size and signedness should be used in place of the basic numerical types. \\
\bottomrule
\end{tabularx}
\end{table}

\begin{important}[Why No Dynamic Memory?]
Rule 21.3 prohibits \texttt{malloc()}/\texttt{free()} because:
\begin{enumerate}[nosep]
  \item Heap fragmentation can cause allocation failures in long-running systems.
  \item Non-deterministic execution time of \texttt{malloc()}.
  \item Memory leaks are catastrophic in systems that cannot be rebooted.
\end{enumerate}
Instead, use statically allocated buffers, memory pools, or stack allocation.
\end{important}

\subsection{CRC for Data Integrity}

Cyclic Redundancy Check (CRC) is widely used for detecting data corruption in communication and storage.

\begin{lstlisting}[style=C, caption={CRC-32 implementation (IEEE 802.3 polynomial).}]
#include <stdint.h>

static uint32_t crc32_table[256];

void crc32_init(void) {
    for (uint32_t i = 0; i < 256; i++) {
        uint32_t crc = i;
        for (int j = 0; j < 8; j++) {
            if (crc & 1)
                crc = (crc >> 1) ^ 0xEDB88320;
            else
                crc >>= 1;
        }
        crc32_table[i] = crc;
    }
}

uint32_t crc32_compute(const uint8_t *data, uint32_t length) {
    uint32_t crc = 0xFFFFFFFF;
    for (uint32_t i = 0; i < length; i++) {
        uint8_t index = (uint8_t)(crc ^ data[i]);
        crc = (crc >> 8) ^ crc32_table[index];
    }
    return crc ^ 0xFFFFFFFF;
}
\end{lstlisting}

% ----------------------------------------------------------------------------
\section{Worked Examples}
\label{sec:ana:examples}

\begin{example}[Reliability Calculation]
A system has three components in series: MCU ($\lambda_1 = 10$ FIT), EEPROM ($\lambda_2 = 5$ FIT), and power supply ($\lambda_3 = 20$ FIT). Calculate: (a) system MTTF, (b) reliability at $t = 10{,}000$ hours, (c) availability if MTTR = 2 hours.

\begin{solution}
(a) Series system failure rate: $\lambda_{\text{sys}} = \lambda_1 + \lambda_2 + \lambda_3 = 35$ FIT.

$\text{MTTF} = \frac{10^9}{\lambda_{\text{sys}}} = \frac{10^9}{35} = 28{,}571{,}429$ hours $\approx 3{,}261$ years.

(b) $R(10{,}000) = e^{-\lambda_{\text{sys}} \times t} = e^{-35 \times 10^{-9} \times 10{,}000} = e^{-3.5 \times 10^{-4}} \approx 0.99965$ ($99.965\%$).

(c) $A = \frac{\text{MTTF}}{\text{MTTF} + \text{MTTR}} = \frac{28{,}571{,}429}{28{,}571{,}429 + 2} \approx 0.99999993$ ($\sim$7 nines).
\end{solution}
\end{example}

\begin{example}[TMR vs.\ Simplex]
Compare the reliability of a simplex system vs.\ TMR at $t = 100{,}000$ hours if each module has $\lambda = 100$ FIT.

\begin{solution}
$\lambda = 100 \times 10^{-9}$ failures/hour. At $t = 100{,}000$ hours:

\textbf{Simplex:} $R_S(t) = e^{-\lambda t} = e^{-100 \times 10^{-9} \times 10^5} = e^{-0.01} = 0.99005$.

\textbf{TMR:} $R_{\text{TMR}}(t) = 3e^{-2\lambda t} - 2e^{-3\lambda t} = 3e^{-0.02} - 2e^{-0.03}$.

$= 3(0.98020) - 2(0.97045) = 2.94060 - 1.94090 = 0.99970$.

TMR reliability $0.99970 > 0.99005$ simplex. TMR improves reliability from 2 nines to nearly 4 nines at this operating point.
\end{solution}
\end{example}

\begin{example}[FMEA Risk Prioritisation]
Given the FMEA table in Table~\ref{tab:ana:fmea}, identify the highest-risk item and propose additional mitigation measures.

\begin{solution}
Highest RPN: \textbf{Firmware stack overflow} (RPN = 224).

This is the highest risk because:
\begin{itemize}[nosep]
  \item Severity is high (8): undefined behaviour can be catastrophic.
  \item Occurrence is moderate (4): stack overflows are common during development.
  \item Detection is poor (7): often manifests as seemingly random crashes.
\end{itemize}

\textbf{Additional mitigations:}
\begin{enumerate}[nosep]
  \item Use MPU (Memory Protection Unit) to mark stack guard regions.
  \item Enable FreeRTOS \texttt{configCHECK\_FOR\_STACK\_OVERFLOW} (method 2).
  \item Run static analysis (e.g., StackAnalyzer) to compute maximum stack depth.
  \item Place stack canary values and check periodically.
  \item During testing, fill stacks with a known pattern and measure high-water mark.
\end{enumerate}

After mitigation, re-rate: Detection drops to 3, new RPN = $8 \times 4 \times 3 = 96$.
\end{solution}
\end{example}

\begin{example}[Formal Verification Property]
Express the following requirements as LTL formulas:
(a) ``The motor is never on while the safety guard is open.''
(b)~``Every sensor reading request is eventually served.''
(c)~``The system never enters both heating and cooling modes simultaneously.''

\begin{solution}
(a) $\square(\text{guard\_open} \rightarrow \neg\text{motor\_on})$

(b) $\square(\text{sensor\_request} \rightarrow \lozenge \text{sensor\_served})$

(c) $\square \neg (\text{heating} \land \text{cooling})$

To verify these in SPIN, they would be written as:
\begin{lstlisting}[language=C, style=C, numbers=none]
/* (a) */ ltl safe_guard { [](guard_open -> !motor_on) }
/* (b) */ ltl liveness   { [](sensor_request -> <>sensor_served) }
/* (c) */ ltl exclusive  { [](!(heating && cooling)) }
\end{lstlisting}
\end{solution}
\end{example}

% ----------------------------------------------------------------------------
\section{Practice Problems}
\label{sec:ana:practice}

\begin{exercise}[Performance Profiling]
You measure a function at 2{,}000 cycles on Cortex-M4 at \freq{168}{MHz}. The function is called every 1\,ms. (a) What is the CPU utilisation due to this function? (b) If you optimise it to 800 cycles using a lookup table that consumes 4\,KB of Flash, discuss the trade-off.
\end{exercise}

\begin{exercise}[Reliability Series-Parallel]
Design a system with two subsystems: Subsystem A ($\lambda_A = 50$ FIT) needs 99.99\% reliability over 5 years. (a) Does a single module meet this? (b) Design a TMR configuration and calculate its reliability. (c) At what operating time does TMR become less reliable than simplex?
\end{exercise}

\begin{exercise}[FMEA Practice]
Create an FMEA table for a drone flight controller with at least 5 failure modes. Include: IMU sensor failure, GPS loss, motor ESC failure, battery voltage drop, and communication link loss. Assign S, O, D ratings and calculate RPNs.
\end{exercise}

\begin{exercise}[Model Checking]
Write a Promela model for two tasks sharing a critical section protected by Peterson's algorithm. Define an LTL property for mutual exclusion and describe how SPIN would verify it.
\end{exercise}

\begin{exercise}[MISRA-C Review]
Identify MISRA-C 2012 violations in the following code and explain how to fix each:
\begin{lstlisting}[style=C, numbers=none]
int* ptr = (int*)malloc(sizeof(int) * 10);
for (int i = 0; i <= 10; i++)  // potential off-by-one
    ptr[i] = i;
printf("Done");               // return value unused
free(ptr);
\end{lstlisting}
\end{exercise}

\begin{exercise}[CRC Application]
A 64-byte data packet is protected by CRC-32. (a) What is the total frame size? (b) If a single bit is flipped in the data, will CRC-32 detect it? (c) What is the probability of an undetected error for a random error pattern?
\end{exercise}
